\section*{Chapter \ref{ch:m3:automated-market-makers} --- Automated Market Makers}

\begin{solution}{exo:m3:automated-market-makers:1}
$3\,000\,000 \times 9.97/(1\,000 + 9.97) = 29\,614.74$ dollars, an average of 2\,961.47.
\end{solution}

\begin{solution}{exo:m3:automated-market-makers:2}
$2\sqrt4/(1 + 4) - 1 = -20\%$.
\end{solution}

\begin{solution}{exo:m3:automated-market-makers:3}
If one coin \omterm{def:m3:blockchains-for-traders:stable}{depegs}, arbitrageurs sell it into the pool at close to one for one until the pool holds mostly that coin: the
providers end up owning the fallen coin, bought near par, as the writer of a put struck near par would.
\end{solution}

\begin{solution}{exo:m3:automated-market-makers:4}
$0.8^2/8 = 8\%$ of pool value a year, about 2.19 basis points a day.
\end{solution}

\begin{solution}{exo:m3:automated-market-makers:5}
LVR: $50\text{ million} \times 0.25/365/8 = \$4\,281$ a day; fees: $20\text{ million} \times 0.05\% = \$10\,000$. Yes: turnover of
40\% against a requirement of 17.1\%.
\end{solution}

\begin{solution}{exo:m3:automated-market-makers:6}
Under the risk-neutral measure the rebalancing portfolio's expected gain is zero, so the expected pool-versus-hold difference equals
the expected LVR. But \omterm{def:m3:automated-market-makers:il}{impermanent loss} also contains the market-risk difference between the pool and holding the initial \omterm{def:m3:blockchains-for-traders:key}{tokens},
which a hedge removes; LVR is what remains after hedging, and it accrues path by path whatever the price does.
\end{solution}

\begin{solution}{exo:m3:automated-market-makers:7}
About 37.0, 37.0 and 37.7 basis points over 30 days with 400 paths, every 15 minutes, hourly and daily. Without fees the expected LVR
does not depend on how often the arbitrage happens: the pool's value $2L\sqrt{P}$ loses $1 - e^{-\sigma^2t/8}$ in expectation whatever
the steps (36.9 basis points here). With fees, less frequent arbitrage lets larger gaps build before trading, and the frequency then
matters.
\end{solution}

\begin{solution}{exo:m3:automated-market-makers:8}
Fees are only the income side. The return also includes the loss to arbitrageurs (LVR) and the market exposure of the \omterm{def:m3:blockchains-for-traders:key}{tokens} held;
net of LVR and hedged, the return can be negative although fees were 12\%.
\end{solution}

\begin{solution}{pb:m3:automated-market-makers:1}
\textbf{1.} $0.6^2/365/8 = 1.23$ basis points a day; USD~1\,233 a day on USD~10 million. \textbf{2.} $6\text{ million} \times 0.30\%
\times 10/200 = \$900$. \textbf{3.} 4.11\% of pool value a day, USD~8.2 million for this pool. \textbf{4.} $(900 - 1\,233) \times 365
\approx -\$121\,500$. \textbf{5.} LVR quadruples: about $-\$1.47$ million a year. \textbf{6.} Worse at the same volume: fees fall to
\$150 a day; a lower fee must attract several times the volume to pay. \textbf{7.} No: while the price is in range, fees and LVR both
scale with the concentration. \textbf{8.} Hedges the ether exposure (short perpetuals or futures sized to the pool's ether holding and
rebalanced), paying funding, fees and the hedge's own rebalancing cost. \textbf{9.} 1.23\% a day, if volume did not fall with the
higher fee. \textbf{10.} It takes the fees of large swaps without bearing their arbitrage, lowering passive providers' income per unit of
LVR. \textbf{11.} Classify swaps: those that move the pool towards the centralised-exchange price within the block are arbitrage; the
rest, from known routers and wallets, are noise. \textbf{12.} Mark the pool's holdings to a centralised-exchange price at each block and
sum the differences between the rebalancing portfolio and the pool, as the paper does. \textbf{13.} Longer blocks let larger mispricings
build before arbitrage; with fees, arbitrage waits for gaps beyond the fee, and the loss per trade depends on the block time.
\textbf{14.} Little: over a month it is dominated by the price change; LVR is visible only after hedging. \textbf{15.} Hedged P\&L, fees,
LVR, volume mix, the range (if concentrated) and gas. \textbf{16.} Arbitrageurs and block builders win; passive providers lose unless
noise volume is high; noise traders pay fees. \textbf{17.} Only partly: a passive provider quotes a fixed curve, cannot move its price
when it learns, and pays the informed. \textbf{18.} Their fees per unit of LVR are high: \omterm{def:m3:blockchains-for-traders:stable}{stablecoin} pools have tiny $\sigma$ and heavy
volume. \textbf{19.} \emph{Named result:} 4.11\% of pool value a day at 60\% volatility and 0.30\%; this pool turns over 3\% and does not
clear it. \textbf{20.} An option to be picked off by anyone who knows the price, paid for with fees.
\end{solution}

\begin{solution}{iq:m3:automated-market-makers:1}
$(x + (1 - \gamma)\Delta x)(y - \Delta y) = xy$, so $\Delta y = y(1 - \gamma)\Delta x/(x + (1 - \gamma)\Delta x)$.
\iqlookfor{only the post-fee input counts towards the invariant.}
\end{solution}

\begin{solution}{iq:m3:automated-market-makers:2}
The pool's value against holding the initial \omterm{def:m3:blockchains-for-traders:key}{tokens} after a price move: $2\sqrt r/(1 + r) - 1$ for a full-range pool. The name suggests it
disappears; it does only if the price returns, and it mixes a hedgeable market exposure with the real cost, which LVR isolates.
\iqlookfor{the formula, and the LVR distinction.}
\end{solution}

\begin{solution}{iq:m3:automated-market-makers:3}
Instantaneous LVR is $\tfrac12\sigma^2P^2|x^{*\prime}(P)|$; with $x^* = L/\sqrt P$ it is $\tfrac14L\sigma^2\sqrt P$, and over the pool value
$2L\sqrt P$ it is $\sigma^2/8$.
\iqlookfor{the rebalancing argument and the constant-product holding.}
\end{solution}

\begin{solution}{iq:m3:automated-market-makers:4}
Contracts cannot use floating point; square-root prices make swap updates linear within a range. Off-chain reimplementations must reproduce
rounding direction (down for outputs, up for inputs), fixed-point formats and tick boundaries exactly, or their quotes and simulated fills will
disagree with the chain by a unit, which is enough to fail a transaction or mis-size an arbitrage.
\iqlookfor{exactness to the last unit.}
\end{solution}

\begin{solution}{iq:m3:automated-market-makers:5}
Compare all-in costs: price impact by pool, fees, gas, and exposure to sandwiching for a public pool swap; an aggregator's split; an intent's
expected fill (auction dynamics, fillers' competitiveness) and its failure risk; and the centralised-exchange alternative. Simulate each with
the current state and the firm's own history of fills.
\iqlookfor{all-in cost including MEV exposure.}
\end{solution}

\begin{solution}{iq:m3:automated-market-makers:6}
Update prices faster than arbitrageurs can trade (an oracle-fed curve), charge dynamic fees that rise with volatility, auction the right to the
first trade of each block and return the proceeds to providers, or batch trades. Each costs something: \omterm{def:m3:blockchains-for-traders:contract}{oracle} risk, higher fees for noise
traders, complexity, or latency.
\iqlookfor{attack the stale price that LVR comes from, and name the cost.}
\end{solution}
